% intro.tex (국문 v10)
활동층은 영구동토 위에서 여름에 녹고 겨울에 다시 어는 지층이다. 활동층 두께(ALT)는 계절 융해의 깊이, 기반 시설의 안정성, 분해에 노출되는 토양 탄소의 양을 정한다\cite{Hjort2018,Schuur2015,Biskaborn2019}. ALT는 여름 열량에 따라 해마다 달라지며, CALM(Circumpolar Active Layer Monitoring) 같은 관측망이 기록하는 것도 이 연도별 변화이다\cite{Nelson1997,Streletskiy2025}. 그러나 관측은 드물고 고르지 않다. 대부분의 측정은 소수의 집중 관측과 장기 격자에서 나오고, 한 번만 방문한 지점이 많다\cite{Moore2025,Veremeeva2025}.

북반구 영구동토 지역에는 1 km ALT 제품이 여러 개 있다. ESA CCI Permafrost 같은 과정 기반 추정은 위성 지표 온도로 열전도 모델을 구동한다\cite{Westermann2024,Westermann2025}. 통계 · 기계학습(ML) 제품은 지점 관측과 기후 · 지형 · 토양 공변량의 관계를 학습한다\cite{Aalto2018,Ran2022a,Liu2024,Wei2026}. 이 제품들은 대개 여러 해의 평균 조건이나 여러 해의 지점 관측을 합쳐 학습하고, 서로 다른 지점 집합과 검증 방식으로 평가된다. 따라서 제품 값을 특정 셀 · 특정 해의 관측과 바로 비교할 수 없고, 같은 셀에서도 제품 간 차이가 수십 cm에 이른다. 검증 설계도 불확실성을 더한다. 공간적으로 모인 관측에 무작위 교차검증을 쓰면 새 장소에 대한 ML 모델의 정확도가 과대평가된다\cite{Roberts2017,Ploton2020,MeyerPebesma2022,Wadoux2021}.

물리 모델은 다른 접근을 제공한다. Stefan 해 $\mathrm{ALT}=E\sqrt{\mathrm{TDD}}$는 융해 깊이를 그 계절의 융해 도일(TDD)과 하나의 계수 $E$로 잇는다. $E$는 토양의 열적 성질과 수분에 따라 정해진다\cite{Stefan1891,Riseborough2008}. 이 식은 계절 융해의 진행을 재현하고 해마다의 기후에 반응하지만, 지역마다 계수 하나를 쓰면 토양 유기탄소와 지형에 따라 체계적인 잔차가 남는다. 물리 모델이 주된 반응을 맡고 ML이 물리 모델이 놓친 부분을 학습하는 결합은 여러 환경 예측에서 정확도를 높였다\cite{Willard2022}. ALT에서도 알래스카를 대상으로 과정 기반 추정과 ML을 결합한 연구가 있다\cite{Gautam2025}. 그러나 입력 기후를 관측 연도와 맞췄을 때, 라벨이 없는 지역으로 옮겼을 때 이런 결합이 어떻게 작동하는지, 새 지역에 라벨이 몇 개 필요한지는 밝혀지지 않았다.

이 연구는 셀-연도 설계로 이 질문들을 다룬다. 모든 ALT 관측은 특정 해의 1 km 셀이고, 모든 기후 입력은 그 해의 ERA5-Land 자료로 관측일까지 누적해 계산한다\cite{munozsabater2021_era5_land}. 알래스카, 레나 델타, 캐나다 북서부의 2010--2024년 셀-연도 관측 \nncyAll{}개로 Stefan 모델, 직접 ML, 재보정한 Stefan 모델의 잔차를 gradient boosting 으로 학습하는 물리 잔차 워크플로를 비교하였다. 모델은 다섯 가지 검증 설계, 지역 안 · 지역 간 라벨 수 사다리, 외부 영구동토 관측 지역, 같은 셀 · 같은 해의 공개 제품 6종과의 비교로 평가하였다. 결과로 (1) 관측 연도 기후 입력의 효과, (2) 물리 잔차 모델의 Stefan 모델 대비 이득과 그 조건, (3) 새 지역에 필요한 라벨 수, (4) 공개 제품과의 차이를 보고하고, 2011--2024년의 연도별 1 km 지도를 제시한다.
